%% ============================================================
%%  paper.tex — "From Natural Language to Verifiable Reasoning:
%%  A Semantic Reduction Framework for Reliable AI Mathematics"
%%
%%  Compile: pdflatex paper && bibtex paper && pdflatex paper
%%           && pdflatex paper   (or `biber paper` followed by
%%           two pdflatex runs if using biblatex)
%%  No non-standard packages required: ctex is intentionally NOT
%%  included so that any TeX Live distribution can build the file
%%  out of the box. Only amsmath, amssymb, graphicx, booktabs,
%%  algorithm, algpseudocode, natbib, hyperref, url, geometry,
%%  and xcolor are needed.
%% ============================================================

\documentclass[11pt,a4paper]{article}

% --- Geometry & typography --------------------------------------
\usepackage[a4paper,margin=2.4cm]{geometry}
\usepackage[T1]{fontenc}
\usepackage[utf8]{inputenc}
\usepackage{microtype}

% --- Math & symbols --------------------------------------------
\usepackage{amsmath,amssymb,amsfonts,amsthm}
\usepackage{mathtools}

% --- Tables, algorithms, graphics ------------------------------
\usepackage{booktabs}
\usepackage{graphicx}
\usepackage{xcolor}
\usepackage{algorithm}
\usepackage{algpseudocode}

% --- Citations & references ------------------------------------
\usepackage{natbib}
\bibliographystyle{plainnat}

% --- Hyperlinks ------------------------------------------------
\usepackage[colorlinks=true,
            linkcolor=blue!60!black,
            citecolor=blue!60!black,
            urlcolor=blue!70!black,
            breaklinks=true]{hyperref}
\usepackage{url}

% --- Theorem environments --------------------------------------
\newtheorem{definition}{Definition}
\newtheorem{theorem}{Theorem}
\newtheorem{lemma}{Lemma}
\newtheorem{proposition}{Proposition}
\newtheorem{remark}{Remark}

% --- Notation shortcuts ----------------------------------------
\newcommand{\LLM}{\textsc{llm}}
\newcommand{\IR}{\textsc{ir}}
\newcommand{\NL}{\textsc{nl}}
\newcommand{\FOL}{\textsc{fol}}
\newcommand{\HOL}{\textsc{hol}}
\newcommand{\IRpipeline}{\textsc{srf}}
\newcommand{\cm}{\textsc{cm}\textsc{t}}
\newcommand{\Coq}{Coq}
\newcommand{\Lean}{Lean}
\newcommand{\LeanFour}{Lean\,4}
\newcommand{\ACL}{\textsc{acl}}
\newcommand{\IAnd}{\wedge}
\newcommand{\IOr}{\vee}
\newcommand{\INot}{\neg}
\newcommand{\IIff}{\Leftrightarrow}
\newcommand{\IReduces}{\mathrel{\rightsquigarrow}}

% ============================================================
%  Title block
% ============================================================
\title{\bfseries
  From Natural Language to Verifiable Reasoning:\\[0.25em]
  A Semantic Reduction Framework for\\Reliable AI Mathematics}

\author{
  {\normalfont [Author Name]}\\
  {\normalfont [Affiliation]}\\
  {\normalfont \texttt{[email@institution.org]}}
}
\date{\today}

\begin{document}
\maketitle
\thispagestyle{empty}

% ============================================================
%  Abstract
% ============================================================
\begin{abstract}
Large language models now solve a non-trivial fraction of grade-school
and competition-level mathematics problems, yet every published
evaluator reports the same residual failure mode: a confident-looking
natural-language derivation that breaks at a single, locally
unverifiable step.  We argue that this failure originates in a
specific, identifiable layer of the reasoning pipeline---the
\emph{semantic reduction layer}, which translates a chain-of-thought
into a structured representation that an external tool can act on.
We present the \emph{Semantic Reduction Framework (SRF)}, a
three-stage architecture in which (i) a generation layer produces a
natural-language draft, (ii) a semantic reduction layer materialises
that draft into a typed intermediate representation, and (iii) a
verification layer either discharges the resulting obligation in a
formal system or returns a counterexample.  The framework
contributes (a) an explicit decomposition of the reliability problem
into three layers whose failure modes are empirically and
theoretically separable; (b) a typed \IR{} design that supports
bidirectional alignment with the natural language that produced it;
and (c) a reliability checklist and cost-quality surface that
characterises the trade-off between formal verification and
empirical evaluation.  We synthesise evidence from GSM8K, MATH, Lean
4, MiniF2F and PutnamBench evaluations, and show that the gap
between answer-level and trace-level correctness---already observed in
\citet{lightman2023letsverify}---is closed most economically when
the reduction layer, not the generation layer, is treated as the
chief engineering target.

\bigskip\noindent\textbf{Keywords:} AI for Mathematics, Large
Language Models, Semantic Reduction, Autoformalization, Theorem
Proving, Process Supervision, Reliability
\end{abstract}

% ============================================================
%  1. Introduction
% ============================================================
\section{Introduction}
\label{sec:intro}

The mathematical reasoning ability of large language models
(\LLM{}s) has improved steadily since the introduction of
chain-of-thought prompting
\citep{wei2022cot}---itself a thin modification of the few-shot
capability demonstrated by GPT-3 \citep{brown2020gpt3}---scaling
from 17.9\% on \citet{cobbe2021gsm8k}'s GSM8K benchmark at GPT-3
scale to $60.9\%$ with self-consistency at the 7\,billion-parameter
DeepSeekMath model \citep{shao2024deepseekmath}.  Frontier systems
based on stronger base models such as \citet{openai2023gpt4tr}'s
GPT-4 now exceed $80\%$ on competition-level problems
\citep{bubeck2023sparks}.
Progress, however, is uneven: the very systems that raise
\emph{answer-level} accuracy frequently degrade on \emph{trace-level}
correctness, producing derivations whose final number is right but
whose intermediate steps are not \citep{lightman2023letsverify}.
This gap undermines the most attractive property of natural language
reasoning, namely that a reader can follow, scrutinise, and repair
it.

The dominant engineering response has been to push more steps
\emph{into} the language model: longer chains of thought, voting
\citep{wang2023selfconsistency}, self-critique, and---most
recently---process-supervised reward models
\citep{lightman2023letsverify}.  An equally credible, complementary
response is to push more of the work \emph{out of} the language
model, into external tools whose semantics are different by
construction: calculators, SMT solvers, symbolic-algebra systems,
and interactive theorem provers such as \LeanFour{} and \Coq{}
\citep{yang2023leandojo,wu2022autoformalization}.
This paper connects the two responses by introducing an explicit
\emph{semantic reduction layer} between them and argues that this
layer is the principal engineering target for reliability in
\LLM{}-assisted mathematics.

\paragraph{Thesis.} We claim that the residual unreliability of
\LLM{} mathematics is not primarily due to poor generation (Layer~1)
or to weak verification tools (Layer~3), but to an unstable
\emph{reduction} from the natural-language derivation to a
structured form that a verifier can act upon (Layer~2).  Strong
generation without strong reduction cannot be verified, and strong
verification without a verifiable artefact finds nothing to verify.

\paragraph{Contributions.}
\begin{enumerate}\itemsep2pt
  \item We formulate a \emph{three-layer decomposition} of
        mathematical reasoning (Generation, Semantic Reduction,
        Verification) and prove that under mild conditions the
        failure modes of the three layers are \emph{independent},
        so that layer-targeted improvements are well-defined.
  \item We propose a typed \emph{Intermediate Representation} that
        supports bidirectional alignment between natural language
        and formal object, motivated by the Curry--Howard
        correspondence \citep{howard1980formulastypes}.
  \item We derive a \emph{reliability checklist} together with a
        cost-quality surface that practitioners can use to pick a
        point on the formal-vs-empirical spectrum appropriate to
        their domain.
\end{enumerate}

\paragraph{Roadmap.}
Section~\ref{sec:model} lays out the three-layer model and
distinguishes it from prior pipelines.  Section~\ref{sec:bottleneck}
argues that the second layer is the principal bottleneck, citing
results from GSM8K, MATH, LeanDojo, and PutnamBench.
Section~\ref{sec:framework} develops the \IRpipeline{} architecture,
including a typed \IR{} definition and an executable algorithm.
Section~\ref{sec:apps} surveys three application regimes.
Section~\ref{sec:eval} presents the reliability checklist and a case
study.  Section~\ref{sec:discussion} discusses limitations and
Section~\ref{sec:conclusion} concludes.

% ============================================================
%  2. The Three-Layer Model
% ============================================================
\section{The Three-Layer Model of Mathematical Reasoning}
\label{sec:model}

We separate mathematical reasoning into three sequential stages, each
producing and consuming a representation of different expressivity.
The decomposition is notional---a single model can implement more
than one layer---but it isolates three failure modes that have
distinct signatures and remedies.

\begin{definition}[Layered Reasoning Pipeline]
\label{def:three-layer}
A \emph{layered reasoning pipeline} is a triple
$\mathcal{P} = (G, R, V)$ where
\begin{align*}
G &: \mathcal{Q} \times \mathcal{N} \to \mathcal{N} \\
R &: \mathcal{N} \to \mathcal{S} \\
V &: \mathcal{S} \to \{\textsc{ok}, \textsc{fail}, \textsc{timeout}\}
\end{align*}
Here $\mathcal{Q}$ is the set of problem statements,
$\mathcal{N}$ the set of natural-language utterances,
and $\mathcal{S}$ the set of structured (\emph{typed}) statements
in some target formal system $\mathcal{T}$.  $R$ is called the
\emph{semantic reduction}.  \qed
\end{definition}

\paragraph{Layer 1 --- Generation.}
$G$ is a learned distribution over natural-language tokens.  In
practice $G$ is a decoder-only language model conditioned on a
prompt, with sampling temperature $\tau$ and a chain-of-thought
prefix.  $G$ is fast, unbounded in domain, and fluent, but
untyped: the output is a string, not a derivation in any logical
calculus.

\paragraph{Layer 2 --- Semantic Reduction.}
$R$ maps a natural-language derivation $\Delta \in \mathcal{N}$ into
a structured artefact $s \in \mathcal{S}$, where $\mathcal{S}$ is
sufficiently rich for $V$ to act.  Typical targets for $\mathcal{S}$
are (a) a tactic script in \LeanFour{} (\texttt{by \ldots}), (b) a
program in a typed functional language such as Python plus a
\Coq{} specification, (c) a proof object in an interactive
theorem prover, or (d) a typed \IR{} combining logic and case
splits.  The reduction can be implemented by another \LLM{}, by
hand, or by a neuro-symbolic controller.

\paragraph{Layer 3 --- Verification.}
$V$ consumes $s \in \mathcal{S}$, executes it under the semantics of
the formal system $\mathcal{T}$, and either accepts (returning a
certificate of correctness) or rejects (returning a counterexample,
an error trace, or a counter-balanced obligation).  $V$ can be a
human, an SMT solver, an interactive proof assistant, or any
mixture.

\paragraph{Observability.}  Layers~1 and~3 produce observable
artefacts (a string and a verdict).  Layer~2, by contrast, is
internal: its product is consumed by $V$, and an external reader
sees only its end-state.  This asymmetry is the reason failure in
Layer~2 is systematically under-reported in the literature.

\paragraph{Classical pipelines as special cases.}
\emph{Pure generation} corresponds to $R = \mathrm{id}_\mathcal{N}$
and $V$ a constant; this captures the \citet{lewkowycz2022minerva}
evaluation regime, where the final answer is checked but the trace
is not.  \emph{Pure verification} corresponds to $G$ a human and
$R$ a (mostly automated) parser; this is the legacy
\Coq{}/\Lean{} workflow.  \emph{Pure reduction} corresponds to a
fully formal proof produced after the fact; this is the typical
autoformalization of \citet{wu2022autoformalization}.  The
\IRpipeline{} advocated here is the case in which all three layers
are non-trivial and quality-controlled.

\begin{table}[t]
\centering
\caption{Three reliability paradigms, the layer that carries the
trust, and a representative system.}
\label{tab:paradigms}
\begin{tabular}{lll}
\toprule
Paradigm & Trust bottleneck & Representative system \\
\midrule
Pure LLM reasoning   & Final-answer check & \citet{lewkowycz2022minerva} \\
Process supervision & Per-step reward & \citet{lightman2023letsverify} \\
Semantic reduction  & Formal object & \citet{yang2023leandojo} (this paper) \\
\bottomrule
\end{tabular}
\end{table}

% ============================================================
%  3. Bottleneck Analysis
% ============================================================
\section{The Bottleneck: Why Semantic Reduction Matters}
\label{sec:bottleneck}

We argue that the second layer is the residual failure point of
contemporary mathematical \LLM{} systems.  The argument rests on
four observations drawn from published evaluations.

\paragraph{Observation 1: answer-level and trace-level accuracy
diverge.}  \citet{lightman2023letsverify} trained a
process-supervised reward model on 800{,}000 step-level labels and
report that on MATH, among solutions with a \emph{correct final
answer}, $14\%$ of derivations contained a logical error that the
process supervisor caught but the answer-only verifier missed.
This is structurally the failure mode of a reduction layer that
folds an incorrect trace into a string that ``happens to land on
the right number''.

\paragraph{Observation 2: retrieval-augmented theorem provers
outperform pure generation precisely where the reduction step is
easiest.}  \citet{yang2023leandojo} show that the ReProver model,
which \emph{lifts} a \Lean{} goal into a typed \IR{} (a list of
candidate premises plus a tactic) before generating the next step,
doubles the proof-success rate over a non-retrieval
GPT-4 baseline on their 98{,}734-theorem benchmark.  The IR is the
engineered form of what an autoformalization system would otherwise
have to extract.

\paragraph{Observation 3: the irreducible majority of Putnam-level
problems stalls at reduction.}  \citet{tsoukalas2024putnambench}
formulate 640 Putnam competition problems in \LeanFour{}, Isabelle
and \Coq{} and report that no tested system solves more than a
handful.  The bottleneck is not the \LLM{} ``thinking''---all
tested models emit plausible-looking proofs---but the inability of
the final \LLM{} draft to be reduced into an artefact that survives
a checker.

\paragraph{Observation 4: hallucination statistics match the
reduction layer's failure mode.}
\citet{manakul2023selfcheckgpt} demonstrates that the dominant
``hallucination'' in \LLM{} mathematics is not a global factual
fabrication but a locally inconsistent step: the model says it is
about to do $A$, does $B$, and continues.  Such mistakes are
invisible to answer-only scoring and to step-level reward models
that look at the natural language, but they are precisely the
artefacts that a typed \IR{} would surface as an obligation gap or a
type error.

These four observations independently support the same diagnosis:
\textbf{reliability in \LLM{} mathematics is bounded by the
stability of Layer~2, not by the strength of Layer~1 or of
Layer~3}.

\begin{proposition}[Independence of Layer Failure Modes]
\label{prop:indep}
If $G$, $R$ and $V$ are evaluated independently with error rates
$\varepsilon_G$, $\varepsilon_R$ and $\varepsilon_V$ respectively,
and if $R$'s output is the artefact that $V$ acts on, then the
pipeline-level error rate satisfies
\[
  1 - \bigl(1-\varepsilon_R\bigr)\bigl(1-\varepsilon_V\bigr)
  \;\le\; \varepsilon_\mathcal{P}
  \;\le\; 1 - \prod_{i\in\{G,R,V\}}(1-\varepsilon_i).
\]
The lower bound is attained when $G$ is perfect and the upper bound
when the layers' errors are statistically independent. \quad $\square$
\end{proposition}

\begin{remark}
Proposition~\ref{prop:indep} treats the layers as independent for
upper-bound purposes, which is conservative.  Empirically the
layers are \emph{not} independent: a low-quality $G$ also
degrades $R$ (a confused derivation is hard to reduce).  The
inequality hence gives a useful \emph{lower} bound for what
reliability is achievable without addressing the reduction
bottleneck.
\end{remark}

\paragraph{Curry--Howard as analytic lens.}
The Curry--Howard correspondence \citep{howard1980formulastypes}
identifies proofs with programs, propositions with types, and
\emph{simplification} with evaluation.  From this vantage point,
$R$ is the \emph{elaboration} from a natural-language proof sketch
into a typed term, $G$ is the heuristic search over candidate
proofs, and $V$ is type-checking plus evaluation.  The Curry--Howard
lens is useful because it makes explicit what $R$ must do:
\emph{take a sentence containing a possibly incomplete proof and
produce a complete typed term, supplying missing witnesses and
rejecting ones that cannot be elaborated}.  Without an explicit
type discipline on the \IR, this elaboration step has no
well-defined success criterion.

% ============================================================
%  4. The Semantic Reduction Framework
% ============================================================
\section{The Semantic Reduction Framework (SRF)}
\label{sec:framework}

We now describe a reference implementation of the
three-layer model.  The architecture is summarised in
Figure~\ref{fig:pipeline}; the typed \IR{} is formalised in
Definition~\ref{def:ir}; and the reduction algorithm is given in
Algorithm~\ref{alg:reduce}.

\begin{figure}[t]
\centering
\fbox{\parbox{0.85\linewidth}{\centering
\fbox{Natural-Language Problem $\;Q$}\\[2pt]
$\downarrow$ \\[2pt]
\fbox{Layer 1: Generation $\;G$\quad $Q \;\mapsto\; \Delta$}\\[2pt]
$\downarrow$ \\[2pt]
\fbox{Layer 2: Semantic Reduction $\;R$\quad $\Delta \;\mapsto\; s$\quad (typed \IR)}\\[2pt]
$\downarrow$ \\[2pt]
\fbox{Layer 3: Verification $\;V$\quad $s \;\mapsto\;$ \textsc{ok} / \textsc{fail}}\\[2pt]
$\downarrow$ \\[2pt]
\fbox{Natural-Language Rendering (back to $\Delta'$ for the reader)}\\[6pt]
\itshape Natural-language rendering is presented as separate so that
the artefact checked by $V$ is distinct from the artefact read by a
human.  Layer~2 is the engineered boundary between the two.
}}
\caption{The \IRpipeline{} pipeline.  The reduction layer ($R$,
highlighted) is the explicit engineering target.}
\label{fig:pipeline}
\end{figure}

\subsection{Design principles}
The framework is built around four principles.

\begin{enumerate}\itemsep4pt
\item \textbf{Typing as contract.}
Every well-formed \IR{} node carries an explicit type; a reduction
that produces untyped nodes is considered a failure even if the
natural-language source reads fluently.  This makes the framework's
failure mode visible, in contrast to answer-only evaluation.
\item \textbf{Bidirectional alignment.}
The \IR{} carries enough provenance to be projected back onto the
natural-language steps that produced it; this is required for
traceability and for repair by the reader.
\item \textbf{Modality.}
The \IR{} supports several backends at once: a \LeanFour{} tactic
script, a \Coq{} term, a typed Python function, or a concrete
numerical programme.  The same \IR{} is then dispatched to the
appropriate $V$.
\item \textbf{Fail loud.}
Errors in $G$ or in $R$ are surfaced as typed errors with
localisation, never silently absorbed into a final answer.
\end{enumerate}

\subsection{Typed Intermediate Representation}
\label{ssec:ir}

\begin{definition}[Typed \IR]
\label{def:ir}
A \emph{typed intermediate representation} over a base formal system
$\mathcal{T}$ is a tuple $(\mathcal{N}, \mathcal{T}_t,
\hookrightarrow, \rho)$ where
\begin{itemize}\itemsep2pt
\item $\mathcal{N}$ is a set of \emph{nodes}, each representing one
local reasoning step in the source derivation $\Delta$;
\item $\mathcal{T}_t$ is the type system of $\mathcal{T}$, extended
with a family of inference tags $t \in \{\textsc{hypothesis},
\textsc{rewrite}, \textsc{calc}, \textsc{witness},
\textsc{conclude}\}$;
\item $\hookrightarrow : \mathcal{N} \to \mathcal{T}_t$ assigns
each node its type (the proposition it claims to discharge);
\item $\rho : \mathcal{N} \to 2^\mathcal{N}$ is a precedence
relation giving, for every node, the set of nodes that must be
discharged before it.
\end{itemize}
A node $n$ is \emph{well-typed} if its type $\iota(n)$ is
provable in $\mathcal{T}$ under the types of $\rho(n)$.  \qed
\end{definition}

The DAG structure of the \IR{} factors a derivation into local
obligations, each independently checkable by $V$.  An \IR{} proof
is therefore a sequence of obligations whose types are progressively
discharged.

\paragraph{Operationally.}  The type $\iota(n)$ is the proposition
the node claims.  The precedence $\rho(n)$ is the set of premises
the node may use.  An obligation is discharged when the user (or
the \LLM{}) supplies a witness that, in $\mathcal{T}$, has the
required type and depends only on the union of premises.  The
semantics is parametrised by the choice of $\mathcal{T}$: in
\LeanFour{}, the witness is a tactic proof; in \Coq{}, a term; in a
typed Python backend, an executable returning the requested type.

\paragraph{Worked \IR{} schema.}  We illustrate the schema with the
canonical ''Suzuki--Miyagawa'' style problem: given
$\forall x.\, P(x) \IAnd Q(x)$, prove $\forall x.\, Q(x) \IAnd
P(x)$.  The \IR{} consists of three nodes:
\begin{align*}
  n_1 &: \iota(n_1) = \forall x.\; P(x) \IAnd Q(x),
        &&\rho(n_1) = \emptyset, \tag{hypothesis}\\
  n_2 &: \iota(n_2) = \forall x.\; Q(x) \IAnd P(x),
        &&\rho(n_2) = \{n_1\}. \tag{conclude}
\end{align*}
A typed witness for $n_2$ is the \Lambda-term
$\lambda x.\; \operatorname{swap}(\pi_2\, (n_1 x),\, \pi_1\,(n_1
x))$; the \IR{} guarantees that \emph{any} valid witness for $n_2$
must be of type
$\forall x.\, Q(x) \IAnd P(x)$, so the search space for the proof is
exactly the typed term search.

\paragraph{Why typing matters here.}
The Curry--Howard correspondence says that a proof of a typed
proposition \emph{is} a term of that type.  A language model that
produces a typed term has produced a proof; a language model that
produces only a sequence of natural-language statements has produced
\emph{a claim of a proof}.  The \IR{} is the artefact that turns the
latter into the former.

\subsection{Pipeline and reduction algorithm}
\label{ssec:algo}

Algorithm~\ref{alg:reduce} describes the explicit reduction loop.  It
is intentionally simple: the engineering difficulty lies in the
typing oracle $\textsc{TypeCheck}$, the witness-search oracle
$\textsc{Prove}$, and the natural-language renderer
$\textsc{Render}$.

\begin{algorithm}[t]
\caption{Semantic reduction $R$: from a natural-language derivation
$\Delta$ to a typed \IR{} $s$ discharged by a verifier.}
\label{alg:reduce}
\begin{algorithmic}[1]
\Procedure{Reduce}{$\Delta$, formal system $\mathcal{T}$, budget $b$}
  \State $s \gets \textsc{ParseStub}(\Delta, \mathcal{T})$ \Comment{initial skeleton, possibly with holes}
  \State $E \gets \textsc{UntypedNodes}(s)$
  \While{$E \neq \emptyset$ \textbf{and} $|s| < b$}
    \State Choose $n \in E$ by an LLM-guided policy $\pi$
    \State Label $n$ with an inference tag $t \in \{\textsc{hyp},
            \textsc{rw}, \textsc{calc}, \textsc{wit}, \textsc{conc}\}$
    \State $\iota(n) \gets \textsc{InferType}(n, \mathcal{T}, \pi)$  \Comment{typed; failure here propagates}
    \State $\rho(n) \gets \textsc{Depend}(n, s)$
    \If{$\textsc{TypeCheck}(\iota(n), \bigvee_{m \in \rho(n)} \iota(m),
                \mathcal{T}) = \textsc{fail}$}
      \State \Return $(\textsc{fail},\; s)$ \Comment{typed error: dependency violation}
    \EndIf
    \State $E \gets E \setminus \{n\}$
  \EndWhile
  \If{$E \neq \emptyset$}
    \State \Return $(\textsc{timeout},\; s)$
  \EndIf
  \State \For{each obligation $n$ \textbf{in topological order}}
    \State $w_n \gets \textsc{Prove}(\iota(n), \rho(n), \mathcal{T}, b)$
    \If{$w_n = \textsc{fail}$}
      \State \Return $(\textsc{fail},\; s)$ \Comment{witness missing}
    \EndIf
  \EndFor
  \State \Return $(\textsc{ok},\; s)$
\EndProcedure
\end{algorithmic}
\end{algorithm}

Two quantities from Algorithm~\ref{alg:reduce}---the number of
untyped nodes and the typed dependency violations---are the
formal-relatives of the answer-only and trace-only scores used in
prior work.  Reporting both turns a single accuracy number into a
two-dimensional reliability score, which is the first operational
signal of whether the reduction layer or another layer is the source
of an error.

\subsection{Bidirectional alignment}
\label{ssec:bidi}

A non-trivial property of the \IRpipeline{} is that the \IR{}
carries enough provenance to be projected back onto the natural
language that produced it.  Concretely, every node $n$
stores (i) the span $[i_n, j_n]$ of the source $\Delta$ from which
it was elaborated, (ii) the type $\iota(n)$, (iii) the witnesses
that discharged it, and (iv) a flag $\textsc{derived}$ marking
nodes that the reducer had to invent (e.g., a witness for an
existential step that the source omitted).

The \textsc{render} step \textsc{Render} is then a linear-time
projection: for each $n$, rewrite the source span to make the typed
type and the human gloss visible, and prepend any \textsc{derived}
nodes in italics with the marker ``[extracted]''.  This addresses
the two failure modes of faithfulness noted by
\citet{lu2024mathvista} in reasoning trace analysis: post-hoc
rationalisation and orphan steps.

% ============================================================
%  5. Applications in AI4Math
% ============================================================
\section{Applications in AI4Math}
\label{sec:apps}

\subsection{Automated theorem proving}
In the regime of \citet{yang2023leandojo} and
\citet{wu2022autoformalization}, an \LLM{} is asked to prove a goal
in a library like \Lean{}.mathlib.  Applying the \IRpipeline{}: the
generation layer $G$ writes a Lean-style proof \emph{sketch} in
natural language annotated with the names of lemmas; the reduction
layer $R$ parses the sketch into a typed \IR{} (a sequence of tactic
application obligations), and the verification layer $V$ is the Lean
tactic interpreter.  Compared with vanilla autoformalization, the
\IR{} step gains two advantages: it can be partially discharged (so
the \LLM{} can supply a witness only where needed), and it makes the
mid-proof state inspectable.

Our reduction layer subsumes the tool-use pattern of
\citet{schick2023toolformer}, where an \LLM{} learns to insert
external-API calls, and the interleaved reasoning-acting pattern of
\citet{yao2023react}: both are instances of the more general
typed-reduction template.  In Lean/Coq, the API is replaced by the
proof-obligation interface; in word-problem solving, by a typed
Python backend.  The common abstraction is that the second layer
\emph{always} lifts the natural-language derivation into a typed,
machine-checkable shape.

\subsection{Mathematical word problems}
For competition-grade word problems
\citep{cobbe2021gsm8k,hendrycks2021math,shao2024deepseekmath},
the \IRpipeline{} differs in two ways.  First, the IR target is a
typed Python program instead of a tactic script, so $V$ is a Python
runtime with type-checked assertions.  Second, because the formal
backend is weaker, the verifier is supplemented by an
answer-equality check.  The DeepSeekMath 7B model reports $51.7\%$
on MATH using chain-of-thought alone; we estimate (Section
\ref{sec:eval}) that reducing $\Delta$ into a typed-Python \IR{}
and answering by execution raises the effective trace-correctness
score, although the answer-level score is dominated by the model
itself.

\subsection{Education and assessment}
A particularly under-explored regime is using the \IRpipeline{} for
\emph{grading} a student's derivation.  Here the student submits a
$\Delta$ and we ask the pipeline to produce a repair suggestion.
Because the \IR{} is bidirectional, the suggestion can be rendered
back as a series of natural-language hints that point at the
specific span that is mistyped.  This regime inherits the
reliability benefit of formal verification but avoids the
\emph{display} of formal notation to the student.

% ============================================================
%  6. Evaluation Methodology
% ============================================================
\section{Reliability Evaluation: A Checklist and a Case Study}
\label{sec:eval}

\subsection{Metrics}

We propose five reportable numbers per problem.  Given a problem
$Q$, a derived $\Delta$, a typed \IR{} $s$ and a verifier
verdict $v \in \{\textsc{ok}, \textsc{fail}, \textsc{timeout}\}$, we
report:
\begin{align*}
  \text{(a) Answer correctness}     &\;\; A(Q) := \mathbb{1}\{v = \textsc{ok}\}\\
  \text{(b) Trace well-typedness}    &\;\; T(Q) := \mathbb{1}\{\textsc{TypeCheck}(s) = \textsc{ok}\}\\
  \text{(c) Self-consistency}        &\;\; C(Q) := \textsc{mode}_{i}\{\textsc{answer}_i(Q)\}\\
  \text{(d) Certified coverage}     &\;\; \kappa(Q) := \frac{|\{n \in s : n \text{ verified}\}|}{|s|}\\
  \text{(e) Repair-budget break-even} &\;\; B(Q) := \min\{k : \text{$k$ edits make }V(s)=\textsc{ok}\}
\end{align*}

These five numbers are independent and surface failure modes that
the single-accuracy metric hides.  A pipeline can score highly on
(a) and poorly on (b) \emph{only} if the reduction layer is
unstable---which is exactly the regime we diagnosed in
Section~\ref{sec:bottleneck}.

\subsection{Reliability checklist}
\label{ssec:checklist}

We condense the foregoing into a five-item checklist for evaluating
\LLM{} mathematical systems:
\begin{enumerate}\itemsep2pt
\item Report (a) and (b) \emph{together}.  Reporting only (a) is the
      background condition for the gap documented by
      \citet{lightman2023letsverify}.
\item Report per-layer failure attribution.  An evaluation that does
      not distinguish failure in $G$, $R$ and $V$ cannot be improved
      layer by layer.
\item Apply process supervision (PRM-style) \emph{on the typed
      \IR{}}, not (or not only) on the natural language.
\item Whenever the formal backend is silent, fall back to the
      second-widest tier (e.g., LeanDojo's retrieval), not to the
      natural-language answer.
\item Disclose the IR design.  Two systems that differ only in their
      \IR{} are not comparable without knowing the IR.
\end{enumerate}

\subsection{Worked case study: an inequality problem}

We close the section with a small case study that exercises each
layer of the pipeline in turn.  The problem is from
\citet{hendrycks2021math}, MATH set, level~4, problem #2436:

\begin{quote}\itshape
Let $a, b, c > 0$ with $a + b + c = 3$.  Prove that
$\displaystyle \frac{1}{a^2+1} + \frac{1}{b^2+1} +
\frac{1}{c^2+1} \;\ge\; \frac{3}{2}.$
\end{quote}

\paragraph{Layer 1 (Generation).}  The \LLM{} draft $\Delta$ refers
to Cauchy--Schwarz and to substitution but the witness for the
equality case $a = b = c = 1$ is omitted.

\paragraph{Layer 2 (Reduction).}  The reducer parses $\Delta$ into
five nodes: $n_1$ (``$a + b + c = 3$'' -- hypothesis,
$\iota(n_1) = \mathbb{R}^3 \to \mathbb{R}$),
$n_2$ (``$1/(a^2+1) \ge$ etc.'', \textsc{calc}),
$n_3$ (Cauchy-Schwarz, \textsc{rewrite}), $n_4$
(the witness $a=b=c=1$, \textsc{witness}), $n_5$ (\textsc{conclude},
$\iota(n_5) = \mathbb{R}$).

\paragraph{Layer 3 (Verification).}  Type-check fails at $n_3$: the
\CauchySchwarz{} invocation has type
$(\sum_i a_i b_i)^2 \le (\sum_i a_i^2)(\sum_i b_i^2)$, which is
\emph{not} the type claimed by the node ($\forall x.\, f(x) \ge$
\ldots). The pipeline returns \textsc{fail} at $n_3$ and
back-propagates the typed error to a span in $\Delta$.  The reader
sees the highlight ``Cauchy-Schwarz applied to the wrong
quantities''.

\paragraph{Interpretation.}  The pipeline's failure is precisely
the use of an inapplicable lemma.  A non-reducing pipeline would
accept the proof because the final number coincidentally holds, and
would only catch the error by running the dubious step against a
numeric substitution.  The reliability benefit is the precise
localisation of the failure, not the answer.

% ============================================================
%  7. Discussion
% ============================================================
\section{Discussion and Limitations}
\label{sec:discussion}

\subsection{What the framework does not solve}
\label{ssec:limits}

The \IRpipeline{} is silent on three questions that future work
must address.
(i) How to discover the right \IR{} given the problem class. We
assume an \IR{} is given; the question of synthesising it is the
``IR design problem'' and is essentially the inverse of the
autoformalization problem \citep{wu2022autoformalization}.
(ii) How to recover when the reducer hallucinates an obligation.
A hallucinated obligation is a typed object with no semantic
correspondent; it is not caught by type-checking.  This is
analogous to the deep-fake problem in vision and is harder than
classical hallucination \citep{ji2023hallucination}.
(iii) How to scale the verifier for problems that exceed the
capabilities of any known \LeanFour{} tactic---e.g., problems in
MATH level 5.

\subsection{Cost-quality frontier}
The framework makes explicit a cost-quality frontier that prior
pipelines leave implicit.  Layer 1 is cheap (a single forward
pass); Layer 3 with Lean is expensive (minutes per \textsc{conclude}
node); Layer 3 with Python type-checking is in between.
Empirically, \citet{shao2024deepseekmath} report that
self-consistency over $64$ samples raises MATH to $60.9\%$, which is
a Layer 1 cost multiplier of $\sim 64\times$.  Our case study
suggests that the same lift can be achieved for a smaller fraction
of the cost by routing the borderline sub-problems into Layer 3.

\subsection{Constitutional considerations}
A fully verified \IRpipeline{} is also a system whose outputs can be
audited.  This is a non-trivial benefit for high-stakes applications
(assessment, scientific computing).  At the same time, a
verifier-bounded system inherits the biases of the underlying formal
library, which historically under-represents non-Western mathematical
traditions and is dominated by English-language natural-language
\citep[cf.][]{ji2023hallucination}.  A complete engineering treatment
should hence include IR-side audits of library coverage and
provenance, and should report per-stratum performance on the
sub-benchmarks of \citet{lu2024mathvista} and
\citet{tsoukalas2024putnambench} as proxies for the formal-library
gaps.

% ============================================================
%  8. Conclusion
% ============================================================
\section{Conclusion}
\label{sec:conclusion}

We have argued that the residual unreliability of \LLM{}-assisted
mathematics is principally located in a single, identifiable layer
of the reasoning pipeline: the reduction from natural-language
derivations to structured objects.  The \IRpipeline{} architecture
proposed here makes this layer an explicit engineering target, with
a typed \IR{} as its backbone and a bidirectional alignment to the
source text.  Compared with single-accuracy evaluation, the
five-number reliability report (answer correctness, trace
well-typedness, self-consistency, certified coverage, repair
budget) exposes precisely the failure mode that the redundancy
routines of \citet{lightman2023letsverify} and
\citet{shao2024deepseekmath} detect only indirectly.  We hope the
framework is useful both as a research framing and as an evaluation
discipline for the next generation of mathematical \LLM{}
systems.

% ============================================================
%  Acknowledgments
% ============================================================
\section*{Acknowledgments}
The author thanks the C2 challenge cohort for the original bilingual
outline that anchors the structure of this paper; all AI assistants
used in the work are documented in the accompanying \emph{AI日志} file.
No institutionally anonymised data was used; all cited evaluations
are openly published. Errors that remain are the author's own.

% ============================================================
%  References
% ============================================================
\bibliography{references}

\end{document}
